\chapter{Fiscal instruments}\label{chap:fiscal}

\section{Revenue-neutral tax replacement}\label{sec:replacement}

The sales tax, import tariff and VAT rates are scaled by $\V{TQADJ}$,
$\V{TMADJ}$ and $\V{TVADJ}$ (in \eqref{eq:P1}, \eqref{eq:P3}, \eqref{eq:P6},
\eqref{eq:P8} and \eqref{eq:I8}), as direct tax rates are scaled by
$\V{TINSADJ}$ \eqref{eq:I10}. All four are fixed at one, except under a closure
that fixes government savings and names one of them as the instrument that
adjusts (\cmd{government(fixgsav direct)}, \cmd{fixgsav sales},
\cmd{fixgsav import} or \cmd{fixgsav vat}; chapter~\ref{chap:closures}).

The instrument scales every base rate of that tax in proportion, so rates that
are zero in the base stay zero; the tax must raise revenue in the base. With
$\V{GADJ}$ fixed, government consumption is unchanged in real terms, so the
closure keeps government savings at their base value: in real terms under the
CPI numeraire, in terms of domestic producer prices under the DPI numeraire.
It replaces lost revenue, for example from a tariff cut, with whichever tax
the user chooses, and the scaling factor shows the rise in rates required.
Under these closures a shock to the instrument's own rates, such as removing
tariffs while tariffs adjust, is not meaningful: the scaling factor offsets
it.

\section{Value added tax with input credits}\label{sec:vat}

VAT is charged on sales of each commodity at the rate $\V{tv}_c$, on the price
$\V{PQ}_c$ before VAT. Final users (households, the government, investment and
the health insurance fund) pay $\V{PH}_c = \V{PQ}_c(1 + \V{tv}_c)$. An activity
recovers the share $\kappa_a$ of the VAT on its inputs as a credit, and so pays
$\V{PQ}_c\,(1 + \V{tv}_c\,(1 - \kappa_a))$ \eqref{eq:P6}. $\kappa_a = 1$, the
default, is a registered activity with full credits; $\kappa_a = 0$ an exempt
or unregistered one, which bears VAT on its inputs. Exports, whether from
domestic output or from composite supply, and purchases of margin services
are zero-rated. Net VAT revenue \eqref{eq:I8} is therefore the VAT on final
demand plus the VAT that activities cannot recover.

\subsection{The SAM}

VAT is recorded in a VAT account whose row collects VAT from the commodity
columns, $\SAM(vat, c)$, net of credits, and whose column pays the government
(\manref{cge_map}, option \cmd{vattax()}). When the SAM records VAT within the
sales tax account, the option \cmd{vatshare()} gives the VAT share of each
commodity's sales tax, which is moved to an internal VAT account before the
checks (\secref{sec:layouts}). VAT cells may be negative.

\subsection{Calibration}

With $\V{PQ}^0_c = 1$, each use $u$ of commodity $c$ recorded in the SAM at the
value $V_{cu}$, at purchaser prices, is a quantity $V_{cu}/(1 + \V{tv}_c\,k_u)$,
where the VAT liability $k_u$ is 1 for final users, $1 - \kappa_a$ for
activity $a$ and 0 for exports and margin services. The rate $\V{tv}_c$ solves
\begin{equation*}
\SAM(vat, c) = \sum_u V_{cu}\,\frac{\V{tv}_c\,k_u}{1 + \V{tv}_c\,k_u},
\end{equation*}
which is increasing in $\V{tv}_c$ when the uses are positive. It is solved by
Newton's method from $\V{tv}_c = \SAM(vat, c) / \sum_u k_u V_{cu}$, to a
relative change of $10^{-15}$; $\V{tv}_c = 0$ when $\SAM(vat, c) = 0$. A
commodity with VAT but no use that bears it is an error. Then
\begin{itemize}
\item $\V{QINT}^0_{ca} = \SAM(c, a) / (1 + \V{tv}_c(1 - \kappa_a))$,
$\V{QINTA}^0_a = \sum_c \V{QINT}^0_{ca}$, $\V{ica}_{ca} = \V{QINT}^0_{ca} / \V{QINTA}^0_a$
and $\V{PINTA}^0_a = \sum_c \SAM(c, a) / \V{QINTA}^0_a$;
\item $\V{QH}^0_{ch}$, $\overline{\V{qg}}_c$, $\overline{\V{qinv}}_c$ and
$\overline{\V{qshi}}_c$ are the SAM values divided by $1 + \V{tv}_c$, and
$\V{PH}^0_c = 1 + \V{tv}_c$;
\item $\V{QQ}^0_c$ (section~\ref{sec:cal-trade}) equals the sum of these
quantities, margin services and exports from composite supply; calibration
checks this.
\end{itemize}
Spending values, budget shares and the accounts of the government and the
fund are unchanged, because they are values at purchaser prices. Without VAT
every quantity equals its SAM value. The VAT rates and credit shares can be
shocked (\cmd{tv}, \cmd{kappa}).

\section{Public debt and interest}\label{sec:debt}

National SAMs record the interest on public debt within government
transfers: to domestic institutions, $\SAM(d, gov)$, and to the rest of the
world, $\SAM(row, gov)$. Given the debt stocks at the start of the base year,
held at home, $\V{DD}^0$ (local currency), and abroad, $\V{DF}^0$ (foreign
currency, equal to local currency at $\V{EXR}^0 = 1$), and their interest
rates $r^d$ and $r^f$, calibration separates the interest from the other
transfers:
\begin{equation*}
\V{trnsfr}_{d,gov} = \SAM(d, gov) - \V{hold}_d\,r^d\,\V{DD}^0,
\qquad
\V{trnsfr}_{row,gov} = \SAM(row, gov) - r^f\,\V{DF}^0,
\end{equation*}
where $\V{hold}_d$ is the share of domestic interest received by institution
$d$. By default it is proportional to $\SAM(d, gov)$, so that interest is the
same share of every institution's transfers; the user may give other shares,
which must be non-negative and sum to one. A non-interest transfer that would
be negative is an error: the SAM's transfers are then too small for the stated
debt and rate. The base equilibrium is unchanged, since at base prices the
interest terms of \eqref{eq:I3}, \eqref{eq:I9} and \eqref{eq:S3} equal the
transfers they replace.

Interest is then paid in nominal terms, on domestic debt in local currency and
on foreign debt in foreign currency, while the other transfers keep their
indexation. A static simulation can change the interest rates or the stocks,
for example a rise in $r^f$ or debt relief that lowers $\V{DF}$ (shocks
\cmd{rd}, \cmd{rf}, \cmd{ddebt}, \cmd{fdebt}).

In a dynamic run the deficit is borrowed, the share $\phi$ abroad
(chapter~\ref{chap:dynamics}); by default $\phi$ is the foreign share of the
base debt, $\V{DF}^0 / (\V{DD}^0 + \V{DF}^0)$. Net borrowing is modelled, not
gross borrowing and amortisation separately: maturing debt is assumed to be
refinanced. The change in the stock of foreign debt measured in local currency
is net foreign borrowing plus a valuation change,
$\V{DF}_{t-1}\,(\V{EXR}_t - \V{EXR}_{t-1})$. Neither default risk nor a link
from debt to interest rates is modelled; the rates are exogenous and can be
changed by shocks.
